\section{Examples}\label{sec:table-examples}

The following examples show what complete tables look like in practice. The first is
a thrust table for a rocket motor.

\lstinputlisting[style=taosmanual]{examples/chapter03/recruits-thr.tbl}

Thrust is a function of time, but the variable \statevar{tmark} is used instead of
\statevar{time} because it can be reset to zero at any point in the trajectory. It is
preferable to \statevar{tseg} because \statevar{tseg} is automatically reset at the
beginning of every segment, whereas \statevar{tmark} is reset only when requested
by the user. The thrust value is multiplied by two factors: the first is the number of
rocket motors, and the second accounts for nozzles canted \ang{6.5} from the vehicle
centerline.

The second example is an axial-force-coefficient table that is only a function of Mach
number. This type of table is commonly used for ballistic sounding rockets that fly
at zero angle of attack.

\lstinputlisting[style=taosmanual]{examples/chapter03/stmi-full.tbl}

The preceding definition uses the full-table format, but it qualifies for the simple
format. In simple form it is

\lstinputlisting[style=taosmanual]{examples/chapter03/stmi-simple.tbl}

The tables shown so far are small, but practical tables are often much larger so that
they more accurately represent a vehicle's characteristics. Trajectory calculations
are often only as accurate as the aerodynamic and propulsion data.

The following mass-flow table represents an Orbus rocket motor. Mass flow is a
function of time, and the data came from an actual motor firing. The table contains
281 paired \statevar{tmark} and \statevar{mdot} values.

\clearpage
\lstinputlisting[
  style=taosmanual,
  basicstyle=\ttfamily\scriptsize
]{examples/chapter03/orbus-mdt.tbl}

The first part of the table contains mass flow in pounds per second as a function of
time. Most of the densely sampled data points are separated by 0.02 seconds. The
second part accounts for mass loss from the attitude-control system.

A final example is an axial-force-coefficient table for a reentry vehicle. It is a
function of Mach number, altitude, and total angle of attack. The original table was
too large to reproduce in its entirety, so the manual gives enough of it to illustrate
the input pattern.

\lstinputlisting[
  style=taosmanual,
  basicstyle=\ttfamily\footnotesize
]{examples/chapter03/reentry-ca-excerpt.tbl}

This table is a function of total angle of attack rather than signed angle of attack
because $C_A$ is symmetric: it has the same value for negative and positive angles
of attack. Using total angle of attack eliminates duplicate tabulated values.

The skewed format is used even though the data satisfies the requirements of the
square format. This is often done when a table is generated by a computer program,
such as Aero,\taosref{20} because the skewed representation is more general and can
be used for any type of tabulated data.
